\section{Introduction}\label{sec:introduction}

Let $M$ and $N$ be von Neumann algebras acting on a complex Hilbert
space $H$, both with identity $I_H$.  Their one-sided near-inclusion gap is
\begin{equation}\label{eq:gamma-definition}
 \gamma(M,N)=\sup_{\substack{x\in M\\\|x\|\leq 1}}
                   \inf_{y\in N}\|x-y\|.
\end{equation}
The approximating element $y$ is not required to be contractive.  A small
value of $\gamma(M,N)$ says that every contraction in $M$ is close to some
element of $N$; it does not require these choices of approximants to
preserve algebraic operations.  If a unitary $u$ satisfies $uMu^*\subseteq N$,
then $\gamma(M,N)\leq 2\|u-I_H\|$.  The converse question asks whether a
small near inclusion can always be obtained in this way.

The underlying perturbation program began with Kadison and Kastler
\cite{KK1972}.  They measured the distance between operator algebras by
the Hausdorff distance between their unit balls in the operator norm,
now called the Kadison--Kastler distance, and proposed that sufficiently
close von Neumann algebras should be related by a unitary close to the
identity.  Christensen, Johnson, and Raeburn--Taylor established the
amenable case in the 1970s
\cite{Chr1977,Joh1977,RT1977}.  Christensen subsequently introduced
near inclusions and proved a spatial embedding theorem for an amenable
source algebra and an arbitrary target
\cite[Definition~2.1 and Theorem~4.3]{Chr1980}.  In particular, sufficiently
small near inclusions in this setting admit implementing unitaries whose
distance from the identity is bounded linearly by the near-inclusion
tolerance.

We use the following formulation of the unrestricted one-sided question
posed by Cameron, Christensen, Sinclair, Smith, White, and Wiggins in
the paragraph preceding \cite[Theorem~2]{CCSSWW2012}:
\begin{quote}
For every $\varepsilon>0$, is there a $\delta>0$ such that, for every
complex Hilbert space $H$ and every pair of von Neumann algebras
$M,N\subseteq\mathcal B(H)$ with common identity $I_H$,
$\gamma(M,N)<\delta$ implies the existence of a unitary
$u\in\mathcal B(H)$ satisfying
\[
 uMu^*\subseteq N
 \qquad\text{and}\qquad
 \|u-I_H\|<\varepsilon?
\]
\end{quote}
Here $\delta$ must be independent of the algebras and their
representations.  We answer this question negatively.

\begin{theorem}\label{thm:main}
There exist $\varepsilon_0>0$, separable complex Hilbert spaces $H_n$,
and von Neumann algebras $M_n,N_n\subseteq\mathcal B(H_n)$ with common
identity $I_{H_n}$ such that
\[
 \gamma(M_n,N_n)\longrightarrow 0,
\]
but, for all sufficiently large $n$, every unitary
$u\in\mathcal B(H_n)$ satisfying $uM_nu^*\subseteq N_n$ obeys
\[
 \|u-I_{H_n}\|\geq\varepsilon_0.
\]
\end{theorem}

In these examples a unitary embedding of $M_n$ into $N_n$ does exist;
the obstruction concerns its distance from the identity.  The theorem
therefore separates normwise approximation from small spatial
implementation even on separable Hilbert spaces.  Its hypothesis is
one-sided.  The construction supplies no reverse near inclusion, and
hence does not give a counterexample to the two-sided Kadison--Kastler
conjecture.  The positive results for nonamenable factors in
\cite{CCSSWW2012}, which concern two-sided perturbations under additional
structural hypotheses, address a different assertion.

The proof transfers a classical finite-dimensional obstruction into the
near-inclusion problem.  Voiculescu's clock and shift matrices
$U_n,V_n\in\mathcal B(\mathbb C^n)$ satisfy
$\|U_nV_n-V_nU_n\|\to0$, but cannot be approximated in norm by commuting
unitary pairs \cite{Voi1983}.  Exel and Loring gave an elementary
winding-number proof of this phenomenon \cite{EL1989}; we use its trace-logarithm form
\cite[Lemma~3.1]{Exel1993}.  We recall the short obstruction argument
in Section~\ref{sec:obstruction}.

To place these matrices in operator algebras, we use the free group
$\F$ on countably many generators.  The starting representation
is the real-parameter Pytlik--Szwarc representation
\cite[\S2.3, Equation~(5)]{PS1986}, whose cyclic coefficient is Haagerup's
positive-definite function $g\mapsto t^{|g|}$, $0<t<1$, where $|g|$ is
reduced word length \cite[Lemma~1.2]{Haa1979}.
Section~\ref{sec:circle} gives a self-contained
tree realization at $t=1/2$ and adds a circle parameter $z$.
The norms of its Fourier coefficients are bounded, uniformly over the
group, by a sequence with summable first moment.  Along successive
distinct free generators, the values of the representations converge
in the weak operator topology to $tz$ times the projection onto the
identity vector in $\ell^2(\F)$.  Substituting $U_n$ and $V_n$ for the
circle parameter therefore produces two unitary representations of
$\F$ whose weak limits retain the original matrices.

Two further steps connect this representation family to the theorem.
In Section~\ref{sec:near}, a diagonal unitary built from the two
representations conjugates a copy of the group von Neumann algebra of
$\F\times\F$.  The target $N_n$ is the commutant of an amplified right
regular representation of this product group.  Reordering the Fourier
expansions gives an approximant in $N_n$ and an error bounded by a
constant times $\|U_nV_n-V_nU_n\|$.  Crucially, this estimate holds on the
entire unit ball of the source algebra.  In
Section~\ref{sec:obstruction}, a hypothetical sequence of spatial
embeddings tending to the identity places the conjugated left regular
representation in $N_n$, where it commutes with the right regular
representation.  Their pointwise product is then a representation of
$\F\times\F$, so the von Neumann algebras generated by its two factors
commute.  Weak operator
limits in these algebras recover operators close to copies of
$U_n$ and $V_n$ supported on projections.  Perturbing those projections,
compressing to a common finite-dimensional corner, and taking polar
parts would yield
commuting unitary matrices close to $U_n,V_n$, contradicting the
winding-number obstruction.
